\documentclass[12pt,letterpaper]{article}
\usepackage{graphicx,textcomp}
\usepackage{natbib}
\usepackage{setspace}
\usepackage{fullpage}
\usepackage{color}
\usepackage[reqno]{amsmath}
\usepackage{amsthm}
\usepackage{fancyvrb}
\usepackage{amssymb,enumerate}
\usepackage[all]{xy}
\usepackage{endnotes}
\usepackage{lscape}
\newtheorem{com}{Comment}
\usepackage{float}
\usepackage{hyperref}
\newtheorem{lem} {Lemma}
\newtheorem{prop}{Proposition}
\newtheorem{thm}{Theorem}
\newtheorem{defn}{Definition}
\newtheorem{cor}{Corollary}
\newtheorem{obs}{Observation}
\usepackage[compact]{titlesec}
\usepackage{dcolumn}
\usepackage{tikz}
\usetikzlibrary{arrows}
\usepackage{multirow}
\usepackage{xcolor}
\newcolumntype{.}{D{.}{.}{-1}}
\newcolumntype{d}[1]{D{.}{.}{#1}}
\definecolor{light-gray}{gray}{0.65}
\usepackage{url}
\usepackage{listings}
\usepackage{color}

\definecolor{codegreen}{rgb}{0,0.6,0}
\definecolor{codegray}{rgb}{0.5,0.5,0.5}
\definecolor{codepurple}{rgb}{0.58,0,0.82}
\definecolor{backcolour}{rgb}{0.95,0.95,0.92}

\lstdefinestyle{mystyle}{
	backgroundcolor=\color{backcolour},   
	commentstyle=\color{codegreen},
	keywordstyle=\color{magenta},
	numberstyle=\tiny\color{codegray},
	stringstyle=\color{codepurple},
	basicstyle=\footnotesize,
	breakatwhitespace=false,         
	breaklines=true,                 
	captionpos=b,                    
	keepspaces=true,                 
	numbers=left,                    
	numbersep=5pt,                  
	showspaces=false,                
	showstringspaces=false,
	showtabs=false,                  
	tabsize=2
}
\lstset{style=mystyle}
\newcommand{\Sref}[1]{Section~\ref{#1}}
\newtheorem{hyp}{Hypothesis}

% highlighting tools to identify final answer
\definecolor{pastelyellow}{RGB}{255, 250, 205}
\usepackage{soul}
\sethlcolor{pastelyellow}

\title{Problem Set 4}
\date{Due: November 26, 2026}
\author{Emily Zhang}


\begin{document}
	\maketitle
	\section*{Instructions}
	\begin{itemize}
		\item Please show your work! You may lose points by simply writing in the answer. If the problem requires you to execute commands in \texttt{R}, please include the code you used to get your answers. Please also include the \texttt{.R} file that contains your code. If you are not sure if work needs to be shown for a particular problem, please ask.
		\item Your homework should be submitted electronically on GitHub.
		\item This problem set is due before 23:59 on Thursday November 26, 2026. No late assignments will be accepted.
	\end{itemize}



	\vspace{.5cm}
\section*{Question 1: Economics}
\vspace{.25cm}
\noindent 	
In this question, use the \texttt{prestige} dataset in the \texttt{car} library. First, run the following commands:

\begin{verbatim}
install.packages(car)
library(car)
data(Prestige)
help(Prestige)
\end{verbatim} 


\noindent We would like to study whether individuals with higher levels of income have more prestigious jobs. Moreover, we would like to study whether professionals have more prestigious jobs than blue and white collar workers.

\newpage
\begin{enumerate}
	
	\item [(a)]
	Create a new variable \texttt{professional} by recoding the variable \texttt{type} so that professionals are coded as $1$, and blue and white collar workers are coded as $0$ (Hint: \texttt{ifelse}).
	
	\lstinputlisting[language=R, firstline=54, lastline=54]{PS04_answersEZ.R}
	
	\textbf{Analysis:}\\
	Here, I created a new column called "professional". \hl{I use the ifelse function to encode the different categories within the type column.} If the type is prof", then it is encoded to 1. Any other type (bc or wc) will be encoded to 0. This esults in professional being a new binary variable. 
	
	\item [(b)]
	Run a linear model with \texttt{prestige} as an outcome and \texttt{income}, \texttt{professional}, and the interaction of the two as predictors (Note: this is a continuous $\times$ dummy interaction.)
	
	\lstinputlisting[language=R, firstline=60, lastline=61]{PS04_answersEZ.R}
	\lstinputlisting[language=R, firstline=64, lastline=66]{PS04_answersEZ.R}
	
% Table created by stargazer v.5.2.3 by Marek Hlavac, Social Policy Institute. E-mail: marek.hlavac at gmail.com
% Date and time: Wed, Nov 25, 2026 - 02:51:32
\begin{table}[H] \centering 
  \caption{Regression of prestige on income and profession} 
  \label{tab:q1} 
\begin{tabular}{@{\extracolsep{5pt}}lc} 
\\[-1.8ex]\hline 
\hline \\[-1.8ex] 
 & \multicolumn{1}{c}{\textit{Dependent variable:}} \\ 
\cline{2-2} 
\\[-1.8ex] & prestige \\ 
\hline \\[-1.8ex] 
 income & 0.003$^{***}$ \\ 
  & (0.0005) \\ 
  & \\ 
 professional & 37.781$^{***}$ \\ 
  & (4.248) \\ 
  & \\ 
 income:professional & $-$0.002$^{***}$ \\ 
  & (0.001) \\ 
  & \\ 
 Constant & 21.142$^{***}$ \\ 
  & (2.804) \\ 
  & \\ 
\hline \\[-1.8ex] 
Observations & 98 \\ 
R$^{2}$ & 0.787 \\ 
Adjusted R$^{2}$ & 0.780 \\ 
Residual Std. Error & 8.012 (df = 94) \\ 
F Statistic & 115.878$^{***}$ (df = 3; 94) \\ 
\hline 
\hline \\[-1.8ex] 
\textit{Note:}  & \multicolumn{1}{r}{$^{*}$p$<$0.1; $^{**}$p$<$0.05; $^{***}$p$<$0.01} \\ 
\end{tabular} 
\end{table}  

	
	\textbf{Analysis:}\\
	I use the lm function in R to run the multi-linear regression with prestige as the outcome and income, professional, and the interaction term income:professional as the predictors. The interaction term is included to see whether the relationship between income and prestige is different for professionals than for blue and white collar workers. From the regression table output, we observe that all the coefficients are signficiant at the alpha threshold of 0.01. This indicates that it is highly unlikely that the relationships observed are due to random chance along. The y-intercept has a value of 21.142, which is the predicted prestige score for a non-professional worker with an income of 0. The standard errors of the coefficients are all relatively small in comparison to their coefficients, so the estimates can be said to be relatively precise.
	
	It is important to note that professional has a larger standard error (4.248) than income (0.0005). This could potentially be explained by the fact that the two variables are measured on different scales. Income is measured in dollars, whereas professional is a binary variable measured by 0 and 1. The standard error may be slightly inflated for professional because the coefficient describes the gap at an income of 0, which is far from the incomes of the professionals in the data, and because professional is strongly correlated with the interaction term. 
	
	Lastly, the fit of the model can be observed through the R-squared value of 0.787, meaning that income, profession type, and their interaction together explain 78.7\% of the variation in prestige. 
	
	\hl{All predictors are signficant at the 0.01 level, and the model explains 78.7\% of the variation in prestige.}
	
	\item [(c)]
	Write the prediction equation based on the result.
	
	\lstinputlisting[language=R, firstline=71, lastline=75]{PS04_answersEZ.R}
	
	\begin{equation*}
		\widehat{prestige} = 21.142 + 0.003 \cdot income + 37.781 \cdot professional - 0.002 \cdot (income \times professional)
	\end{equation*}
	
	\textbf{Analysis:}\\
	To write out the prediction equation, I use the built0in coef( ) function in R to pull the intercept and each slope coefficient. Following the general form of a linear regression with an interaction term, the predicted prestige is equal to the intercept plus each coefficient multiplied by its predictor. It is important to note that the interaction term is the product of income and profesional, so it only contributes to the prediction when professional = 1. 
	
	\item [(d)]
	Interpret the coefficient for \texttt{income}.
	
	\textbf{Analysis:}\\
	Because the model includes an interaction term, the effect of income depends on the value of professional, so the two groups need to be looked at separately. When professional = 0, the prediction equation simplifies to $\widehat{prestige} = 21.142 + 0.003 \cdot income$. When professional = 1, it simplifies to $\widehat{prestige} = 58.923 + 0.001 \cdot income$, since $21.142 + 37.781 = 58.923$ and $0.003 + (-0.002) = 0.001$. This tells us that for those who work blue or white collared jobs, a one unit increase in income is, on average, associated with a 0.003 unit increase in prestige. For those who work professional jobs, a one unit increase in income is, on average, associated with a 0.001 unit increase in prestige. The coefficient of the interaction term (income:professional) shows the difference in slope between these two lines ($0.001 - 0.003 = -0.002$). With that in mind, the average effect of income on prestige does depend on the type of profession.
	
	\hl{The income coefficient (0.003) is the effect of a one unite increase in income on prestige for non-professionals. Prestige tends to be higher for individuals with a higher income, but it increases at a lower rate for professionals than for blue and white collar workers.}
	
	\item [(e)]
	Interpret the coefficient for \texttt{professional}.
	
	\textbf{Analysis:}\\
	The coefficient for professional tells us how much the intercept shifts when we move from the non-professional group to the professional group. When income is set to 0, the interaction term also drops out, so the prediction equation becomes $\widehat{prestige} = 21.142 + 37.781 \cdot professional$. This means that the professional line starts 37.781 units higher than the non-professional line, which we also saw in part d where the professional intercept was 58.923 compared to 21.142 for non-professionals. However, it is important to not that because the model includes the interaction term, the difference between the two groups is not the same at every income level. The gap in predicted prestige is equal to $37.781 - 0.002 \cdot income$, so it shrinks by 0.002 units for every one dollar increase in income. Since no occupation in the data actually has an income of 0, the coefficient is best understood as the difference between the intercepts of the professional and non-professional regression lines. 
	
	\hl{When income is 0, those who have professional jobs have, on average, 37.781 units of prestige more than those who work blue and white collared jobs, and this gap gets smaller as income increases.}
		
	\item [(f)]
	What is the effect of a \$1,000 increase in income on prestige score for professional occupations? In other words, we are interested in the marginal effect of income when the variable \texttt{professional} takes the value of $1$. Calculate the change in $\hat{y}$ associated with a \$1,000 increase in income based on your answer for (c).
	
	\lstinputlisting[language=R, firstline=102, lastline=103]{PS04_answersEZ.R}
	\lstinputlisting[language=R, firstline=106, lastline=106]{PS04_answersEZ.R}
	
	\textbf{Analysis:}\\
	To find the marginal effect of a \$1,000 increase in income for professionals, we use the prediction equation from part c and set professional = 1. We first calculate the predicted prestige when income is 0, and then the predicted prestige when incomeis 1,000. Since the model is linear, the starting value of income does not matter, so an two incomes that are 1,000 apart would give us the same difference. Subtracting the two predictions leaves us with $1000 \times (0.003 + (-0.002))$, which is the slope for professionals multiplied by 1,000. We result in a change of 0.845 in predicted prestige. 
	
	\hl{For professional occupations, a \$1,000 increase in income is associated with a 0.845 unit increase in predicted prestige. }
	
	\item [(g)]
	What is the effect of changing one's occupations from non-professional to professional when her income is \$6,000? We are interested in the marginal effect of professional jobs when the variable \texttt{income} takes the value of $6,000$. Calculate the change in $\hat{y}$ based on your answer for (c).
	
	\lstinputlisting[language=R, firstline=112, lastline=113]{PS04_answersEZ.R}
	\lstinputlisting[language=R, firstline=116, lastline=116]{PS04_answersEZ.R}
	
	\textbf{Analysis:}\\
	To find the marginal effect of changing from a non-professional to a professional occupation at an income of \$6,000, we use the prediction equation and set inceom = 6,000. We calculate the predicted prestige once with professional = 1 and once with professional = 0, and then take the difference. Since income is held at 6,000 in both, the income term cancels out, which leaves us with $37.781 + (-0.002) \times 6000$. This shows that the effect of being a professional gets smaller as inceom increases, because the negative interaction term is subtracted for every unit of income. We result in a change of 23.827 in predicted prestige.
	
	\hl{At an income of \$6,000, changing from a non-professional to a professional occupation is associated with a 23.827 unit change in predicted prestige.}
	
	
\end{enumerate}

\newpage

\section*{Question 2: Political Science}
\vspace{.25cm}
\noindent 	Researchers are interested in learning the effect of all of those yard signs on voting preferences.\footnote{Donald P. Green, Jonathan	S. Krasno, Alexander Coppock, Benjamin D. Farrer,	Brandon Lenoir, Joshua N. Zingher. 2016. ``The effects of lawn signs on vote outcomes: Results from four randomized field experiments.'' Electoral Studies 41: 143-150. } Working with a campaign in Fairfax County, Virginia, 131 precincts were randomly divided into a treatment and control group. In 30 precincts, signs were posted around the precinct that read, ``For Sale: Terry McAuliffe. Don't Sellout Virgina on November 5.'' \\

Below is the result of a regression with two variables and a constant.  The dependent variable is the proportion of the vote that went to McAuliff's opponent Ken Cuccinelli. The first variable indicates whether a precinct was randomly assigned to have the sign against McAuliffe posted. The second variable indicates
a precinct that was adjacent to a precinct in the treatment group (since people in those precincts might be exposed to the signs).  \\

\vspace{.5cm}
\begin{table}[!htbp]
	\centering 
	\textbf{Impact of lawn signs on vote share}\\
	\begin{tabular}{@{\extracolsep{5pt}}lccc} 
		\\[-1.8ex] 
		\hline \\[-1.8ex]
		Precinct assigned lawn signs  (n=30)  & 0.042\\
		& (0.016) \\
		Precinct adjacent to lawn signs (n=76) & 0.042 \\
		&  (0.013) \\
		Constant  & 0.302\\
		& (0.011)
		\\
		\hline \\
	\end{tabular}\\
	\footnotesize{\textit{Notes:} $R^2$=0.094, N=131}
\end{table}

\vspace{.5cm}
\begin{enumerate}
	\item [(a)] Use the results from a linear regression to determine whether having these yard signs in a precinct affects vote share (e.g., conduct a hypothesis test with $\alpha = .05$).
	
	\lstinputlisting[language=R, firstline=122, lastline=122]{PS04_answersEZ.R}
		
	\textbf{Assumptions:}\\
	
	\textbf{Hypothesis:}\\
	
	\textbf{Test Statistic:}
	\lstinputlisting[language=R, firstline=137, lastline=138]{PS04_answersEZ.R}
	
	\textbf{P-Value:}
	\lstinputlisting[language=R, firstline=141, lastline=141]{PS04_answersEZ.R}
	
	\textbf{Conclusion:}\\
	
	\hl{}
	
	\item [(b)]  Use the results to determine whether being
	next to precincts with these yard signs affects vote
	share (e.g., conduct a hypothesis test with $\alpha = .05$).
	
	\textbf{Assumptions:}\\
	
	\textbf{Hypothesis:}\\
	
	\textbf{Test Statistic:}
	\lstinputlisting[language=R, firstline=161, lastline=162]{PS04_answersEZ.R}
	
	\textbf{P-Value:}
	\lstinputlisting[language=R, firstline=165, lastline=165]{PS04_answersEZ.R}
	
	\textbf{Conclusion:}\\
	
	\hl{}
	
	\item [(c)] Interpret the coefficient for the constant term substantively.
	
	\lstinputlisting[language=R, firstline=173, lastline=174]{PS04_answersEZ.R}
	
	\textbf{Analysis:}\\
	
	\hl{}
	
	\item [(d)] Evaluate the model fit for this regression.  What does this	tell us about the importance of yard signs versus other factors that are not modeled?
	
	\lstinputlisting[language=R, firstline=183, lastline=184]{PS04_answersEZ.R}
	
	\textbf{Analysis:}\\
	
	\hl{}
	
\end{enumerate}  


\end{document}
